% ============================================================================
% Section A: Define the System-of-Interest (Deliverable A)
% ============================================================================
\section{Definition of the System-of-Interest}
\label{sec:A}

\subsection{System Identification and Characterization}

\begin{enumerate}
    \item \textbf{Assigned Real-Life System (Annexure A, Group 27):} Electronic-waste-management system (collection, classification, recovery, recycling, hazardous handling, and disposal).
    \item \textbf{Selected System-of-Interest (SoI):} \textbf{\systemOfInterest} --- a formal regional processing plant in MIHAN/Butibori MIDC, Nagpur, rated at \textbf{10 metric tonnes/day (3,000 t/year across 300 days)}.
    \item \textbf{Need or Deficiency:} Over 75\% of urban e-waste leaks into informal scrap channels involving open-air PCB acid leaching and wire burning, causing toxic lead/mercury release, lost strategic metals (Cu, Au, Pd, Al), and producer non-compliance under India's \textit{E-Waste (Management) Rules, 2022}.
    \item \textbf{Principal Purpose:} To safely collect, digitally trace, de-pollute, dismantle, and mechanically recover high-purity secondary raw materials while manifesting hazardous residues to authorized disposal, delivering verifiable Extended Producer Responsibility (EPR) compliance.
    \item \textbf{System Boundary:} Spans from \textit{authorized intake} (12 smart kiosks, 4 GPS vans, weighbridge) through \textit{sorting, manual de-pollution, modular mechanical separation, HEPA/ZLD containment, and EPR logging} up to \textit{verified outbound dispatch} of recovered metals/plastics, PCB lots, and Form-10 hazardous packages.
    \item \textbf{Important External Systems:} Urban generators (households, IT parks, VNIT, hospitals); informal aggregators (\textit{kabadiwalas}); regulators (CPCB EPR Portal, MPCB, NMC, DISH); electronics OEMs/PROs; MSEDCL grid \& MIDC utilities; commodity markets; authorized precious-metal smelters; Butibori Common TSDF; and machinery OEMs.
    \item \textbf{Principal Stakeholders:} Promoter/acquirer SPV, EPC contractor \& equipment OEMs, plant operators, maintenance \& EHS staff, urban generators, electronics OEMs/PROs, regulators (CPCB, MPCB, NMC, DISH), smelters, buyers, TSDF, and local community.
    \item \textbf{Expected Operational Life:} \textbf{15 years} for civil infrastructure, with \textbf{5--7 year refresh cycles} for sorting sensors, EPR software, and high-wear shredder assemblies.
\end{enumerate}

% ============================================================================
% Figure 1: System-of-Interest Context Diagram (Deliverable A)
% Styled in edit.pdf's classic dashed-boundary, periwinkle/rose/orange format
% ============================================================================
\begin{figure}[htbp]
\centering
\resizebox{\textwidth}{!}{%
\begin{tikzpicture}[node distance=1.0cm]

% --- Central Dashed System Boundary (exact edit.pdf aesthetic) ---
\node[rectangle, rounded corners=2pt, draw=black!80, dashed, line width=1.0pt, fill=white,
      minimum width=8.6cm, minimum height=4.4cm] (soi) at (0,0) {};

% Internal sub-blocks inside SoI (soft periwinkle-blue rounded boxes)
\node[sebox, text width=3.55cm, minimum height=0.95cm] (s1) at (-2.05, 0.95)
    {\textbf{1. Collection \& Intake}\\Smart kiosks, vans, QR/RFID};
\node[sebox, text width=3.55cm, minimum height=0.95cm] (s2) at (2.05, 0.95)
    {\textbf{2. Triage \& De-pollution}\\Reuse screen, battery/Hg removal};
\node[sebox, text width=3.55cm, minimum height=0.95cm] (s3) at (-2.05, -0.75)
    {\textbf{3. Mechanical Recovery}\\Shred, eddy/mag/optical sort};
\node[sebox, text width=3.55cm, minimum height=0.95cm] (s4) at (2.05, -0.75)
    {\textbf{4. Containment \& EHS}\\HEPA/ZLD, haz-pack, EPR log};

\draw[searrow] (s1.east) -- (s2.west);
\draw[searrow] (s2.south) -- ++(0,-0.28) -| (s3.north);
\draw[searrow] (s3.east) -- (s4.west);

\node[font=\rmfamily\tiny\itshape\color{black!70}, anchor=south] at ($(soi.south)+(0,0.12)$)
    {Enclosed 4,500 m$^2$ MIHAN/Butibori Facility + 12 Urban Drop-off Kiosks + 4 Reverse-Logistics Vans};

% Bold Boundary Caption underneath (matching edit.pdf, centered cleanly between bottom arrows)
\node[font=\rmfamily\scriptsize\bfseries\color{black}, anchor=north] at ($(soi.south)+(0,-0.10)$)
    {UA-EWRRF System Boundary};

% --- External Systems: Soft Rose Ellipses (matching edit.pdf Pages 5 & 12) ---
\node[seboxExt, text width=2.75cm] (ext_gen) at (-8.3, 1.75)
    {\textbf{Urban Generators}\\Households, IT Parks,\\Hospitals, Inst.};

\node[seboxExt, text width=2.75cm] (ext_inf) at (-8.3, 0.0)
    {\textbf{Informal Sector}\\Kabadiwalas \& Scrap\\Aggregators};

\node[seboxExt, text width=2.75cm] (ext_util) at (-8.3, -1.75)
    {\textbf{Municipal Utilities}\\MSEDCL Power Grid,\\MIDC Water/Telecom};

\node[seboxExt, text width=3.05cm] (ext_reg) at (-2.5, 3.85)
    {\textbf{Regulatory Ecosystem}\\CPCB EPR Portal,\\MPCB, NMC, DISH};

\node[seboxExt, text width=3.05cm] (ext_oem) at (2.5, 3.85)
    {\textbf{Electronics OEMs/PROs}\\Producers Buying\\EPR Credits};

\node[seboxExt, text width=2.85cm] (ext_mkt) at (8.3, 1.75)
    {\textbf{Commodity Market}\\Ferrous, Cu, Al,\\Recycled Polymers};

\node[seboxExt, text width=2.85cm] (ext_smelt) at (8.3, 0.0)
    {\textbf{Authorized Smelters}\\High-Grade PCB Lots\\(Au, Ag, Pd, Cu)};

\node[seboxExt, text width=2.85cm] (ext_tsdf) at (8.3, -1.75)
    {\textbf{Common Haz.\ TSDF}\\Butibori Incineration\\\& Secured Landfill};

\node[seboxExt, text width=3.15cm] (ext_sup) at (-3.35, -4.05)
    {\textbf{Equip.\ OEMs \& Spares}\\Machinery Vendors \&\\Maint.\ Contractors};

\node[seboxExt, text width=3.15cm] (ext_com) at (3.35, -4.05)
    {\textbf{Community \& Env.}\\Ambient Air/Soil/Water\\\& Local Workforce};

% --- Warm Orange Boundary-Crossing Arrows (matching edit.pdf) ---
\draw[seextarrow] (ext_gen.east) -- node[selabel, above, yshift=1pt] {Discarded WEEE} node[selabel, below, yshift=-1pt] {Receipts / Incentives} ($(soi.west)+(0,1.75)$);
\draw[seextarrow, <->] (ext_inf.east) -- node[selabel, above, yshift=1pt] {Bulk Buy-back} node[selabel, below, yshift=-1pt] {Formalized Inflow} (soi.west);
\draw[seextarrow] (ext_util.east) -- node[selabel, above, yshift=1pt] {3-Phase Power / Water} node[selabel, below, yshift=-1pt] {Data Connectivity} ($(soi.west)+(0,-1.75)$);

\draw[seextarrow, <->] (ext_reg.south) -- node[selabel, left, xshift=-2pt] {CTO / Audits} node[selabel, right, xshift=2pt] {EPR \& Manifests} ($(soi.north)+(-2.5,0)$);
\draw[seextarrow, <->] (ext_oem.south) -- node[selabel, left, xshift=-2pt] {EPR Credits} node[selabel, right, xshift=2pt] {Fees / Spec Data} ($(soi.north)+(2.5,0)$);

\draw[seextarrow] ($(soi.east)+(0,1.75)$) -- node[selabel, above, yshift=1pt] {Sorted Metals/Plastics} node[selabel, below, yshift=-1pt] {Commodity Revenue} (ext_mkt.west);
\draw[seextarrow] (soi.east) -- node[selabel, above, yshift=1pt] {Depopulated PCB Batch} node[selabel, below, yshift=-1pt] {Assay Settlement} (ext_smelt.west);
\draw[seextarrow] ($(soi.east)+(0,-1.75)$) -- node[selabel, above, yshift=1pt] {Stabilized Haz-Residues} node[selabel, below, yshift=-1pt] {Form-10 \& Gate Fees} (ext_tsdf.west);

\draw[seextarrow, <->] (ext_sup.north) -- node[selabel, left, xshift=-2pt] {Spares / Blades} node[selabel, right, xshift=2pt] {Failure Logs} ($(soi.south)+(-3.35,0)$);
\draw[seextarrow, <->] (ext_com.north) -- node[selabel, left, xshift=-2pt] {Filtered Air/Jobs} node[selabel, right, xshift=2pt] {Safety Feedback} ($(soi.south)+(3.35,0)$);

\end{tikzpicture}%
}
\caption{Context and System Boundary Diagram for the Urban Automated E-Waste Collection, Dismantling, and Resource Recovery Facility (UA-EWRRF).}
\label{fig:context}
\end{figure}


\subsection{Justification of System Boundary}
\begin{itemize}
    \item \textbf{Included inside the boundary:} Kiosks, collection vans, QR/RFID intake, reuse triage, manual hazardous de-pollution (batteries, mercury lamps, capacitors), modular mechanical recovery, HEPA/carbon filtration, zero-liquid-discharge (ZLD) treatment, and temporary hazardous containment. These are directly engineered by the facility designer and govern purity, safety, and EPR compliance.
    \item \textbf{Kept outside the boundary:} Upstream OEM product design, high-temperature pyrometallurgical smelting (viable only at $>30{,}000$ t/yr), final hazardous incineration/landfilling (mandated to authorized Common TSDFs such as Butibori), municipal landfills, and informal street scrap markets. Keeping these external keeps the SoI tractable while defining strict handover interfaces.
\end{itemize}
